\section{Requirements and Message Protocol}
\label{sec:protocol}

\begin{guidance}[Requirements and message protocol (15 marks)]
Describe and diagram your \textbf{final} protocol: physical link (interface,
baud rate, frame format), \textbf{framing} (terminator, maximum frame length,
timeout), \textbf{board encoding}, the complete \textbf{command set}, and
\textbf{every response and error}. Cover new game, move, win/draw and
disconnect. Include at least one \textbf{sequence diagram}. \textbf{Highlight
every departure} from the brief's reference protocol (p.~3) and justify each
one. A reader should be able to write a compatible client from this section
alone.
\end{guidance}

\subsection{\ph{Heading that states your framing rule, e.g. ``Lines of at most N bytes frame the stream''}}
\ph{Physical link: which USART, which pins, baud rate, data bits, parity, stop
bits, flow control. Framing: the terminator, how you treat \texttt{\textbackslash r},
the maximum frame length and what happens when it is exceeded, the timeout for
an incomplete frame, the character set allowed, the checksum or CRC (if any)
and how it is computed. Use Table~\ref{tab:framing}.}

\begin{table}[t]
\centering
\caption{\ph{Framing rules and the response to each violation.}}
\label{tab:framing}
\footnotesize
\begin{tabular}{@{}L{3.0cm}L{5.0cm}@{}}
\toprule
Rule & Server behaviour \\
\midrule
Terminator & \ph{e.g. \texttt{\textbackslash n} ends a frame} \\
Maximum frame length & \ph{N bytes; what happens on overflow} \\
Incomplete frame & \ph{timeout value; reply sent} \\
Invalid character & \ph{\ldots} \\
Checksum (optional) & \ph{algorithm; reply on mismatch} \\
\bottomrule
\end{tabular}
\end{table}

\subsection{\ph{Heading about your board encoding}}
\ph{How the game state is encoded in a message (e.g. nine characters, cells
0--8 row by row, \texttt{.}/\texttt{X}/\texttt{O}). Give one worked example.}

\subsection{\ph{Heading about your command set}}
\ph{Introduce Table~\ref{tab:commands}. Explain any command that is not obvious,
and how the client learns about changes it did not request (e.g. a game
started from the board's button).}

\begin{table}[t]
\centering
\caption{\ph{Command set. Mark departures from the brief, e.g. with $^\dagger$.}}
\label{tab:commands}
\scriptsize
\begin{tabular}{@{}L{1.9cm}L{1.8cm}L{4.3cm}@{}}
\toprule
Request & Accepted in & Success response \\
\midrule
\texttt{CONNECT} & \ph{state(s)} & \ph{response} \\
\texttt{NEW} & \ph{\ldots} & \ph{\ldots} \\
\texttt{MOVE:<i>} & \ph{\ldots} & \ph{\ldots} \\
\texttt{STATE} & \ph{\ldots} & \ph{\ldots} \\
\texttt{RESIGN} & \ph{\ldots} & \ph{\ldots} \\
\texttt{DISCONNECT} & \ph{\ldots} & \ph{\ldots} \\
\ph{your additions} & \ph{\ldots} & \ph{\ldots} \\
\bottomrule
\end{tabular}
\end{table}

\subsection{\ph{Heading about your errors}}
\ph{List every error response in Table~\ref{tab:errors}, when it is sent and
what it guarantees about the server state. If a request can fail more than one
check, state the order in which you check, so that every request has exactly
one correct response.}

\begin{table}[t]
\centering
\caption{\ph{Every error response and when the server sends it.}}
\label{tab:errors}
\footnotesize
\begin{tabular}{@{}L{2.25cm}L{5.75cm}@{}}
\toprule
Response & Sent when \\
\midrule
\texttt{NACK:BADCMD} & \ph{\ldots} \\
\texttt{NACK:NOGAME} & \ph{\ldots} \\
\texttt{NACK:RANGE} & \ph{\ldots} \\
\texttt{NACK:OCCUPIED} & \ph{\ldots} \\
\texttt{NACK:NOTURN} & \ph{\ldots} \\
\ph{your additions} & \ph{\ldots} \\
\bottomrule
\end{tabular}
\end{table}

\subsection{\ph{Heading about your departures from the brief}}
\ph{For each change or addition to the reference protocol: what you changed
and why it was needed. One or two sentences each. Note: with a single client,
an out-of-turn move may be impossible to send; explain how your protocol makes
\texttt{NACK:NOTURN} testable, or why it does not apply.}

\begin{figure}[t]
\centering
\figph{5.5cm}{Sequence diagram: CONNECT $\rightarrow$ NEW $\rightarrow$ moves
$\rightarrow$ WIN/DRAW $\rightarrow$ STATE $\rightarrow$ DISCONNECT, and/or a
rejected move and recovery from a bad frame. Client on the left, server on the
right, time downwards, replies as dashed arrows.}
\caption{\ph{Caption that states what the diagram shows, e.g. ``A rejected move
leaves the board unchanged.''}}
\label{fig:sequence}
\end{figure}
